% =====================================================================
%  1bat-ud4-13.tex  —  SESSIÓ 13: Buidatge i gràfics amb el full de càlcul
%  (sense preàmbul: s'inclou des de main.tex)
% =====================================================================
\horatitol{Sessió 13 — Buidatge i gràfics amb el full de càlcul}%
{Tancar el cicle d'un projecte de dades: tabular les respostes, treure'n percentatges i generar un gràfic de sectors.}

\subsection{Idea clau}

A la sessió 12 vau \textbf{dissenyar} l'enquesta. Ara en fem el \textbf{buidatge}:
portem les respostes al full de càlcul, les \textbf{tabulem}, en treiem els
\textbf{percentatges} i generem un \textbf{gràfic}. Així tanquem el cicle complet d'un
projecte de dades, exactament el flux que faràs als cicles formatius: dissenyar,
recollir, tabular, calcular i presentar.

\begin{idea}
\textbf{Del full de respostes al gràfic}
\begin{enumerate}[leftmargin=1.6em, itemsep=2pt]
  \item \textbf{Organitza} les respostes en una taula (una fila per persona, o un
        recompte per opció).
  \item \textbf{Compta} quantes vegades surt cada opció amb \texttt{=COMPTA.SI(rang;\,"opció")}
        (\texttt{=COUNTIF}).
  \item \textbf{Percentatge} de cada opció: $\dfrac{\text{recompte}}{N}\per 100$.
  \item \textbf{Gràfic:} selecciona la taula de recomptes i insereix un \textbf{gràfic
        de sectors} (o de barres).
\end{enumerate}
\textbf{Comprova} sempre que els recomptes sumen $N$ i els percentatges, $100\%$.
\end{idea}

\subsection{Exemples}

\begin{exemple}[tabular les respostes d'una pregunta]
A la pregunta «Com de segur et sents a l'institut?» han respost $25$ alumnes. Amb
\texttt{=COMPTA.SI} comptem cada opció i en traiem el percentatge:

\begin{center}
\begin{tabular}{l c c c}
\toprule
\textbf{Resposta} & \textbf{Recompte} & \textbf{\%} & \textbf{Graus (sector)} \\
\midrule
Gens & $2$ & $8\%$ & $29^\circ$ \\
Poc & $5$ & $20\%$ & $72^\circ$ \\
Bastant & $12$ & $48\%$ & $173^\circ$ \\
Molt & $6$ & $24\%$ & $86^\circ$ \\
\midrule
\textbf{Total} & $25$ & $100\%$ & $360^\circ$ \\
\bottomrule
\end{tabular}
\end{center}

\textbf{Comprovació:} els recomptes sumen $25$ i els percentatges, $100\%$.
\end{exemple}

\begin{exemple}[el gràfic de sectors automàtic]
Amb la taula anterior, el full de càlcul genera aquest diagrama de sectors (aquí el
reproduïm). La majoria se sent «Bastant» segura ($48\%$):

\begin{center}
\begin{tikzpicture}[scale=1.5]
  % Gens 0-29
  \fill[udgrey!55] (0,0) -- (0:1) arc (0:29:1) -- cycle;
  % Poc 29-101
  \fill[udblue!30] (0,0) -- (29:1) arc (29:101:1) -- cycle;
  % Bastant 101-274
  \fill[udblue!70] (0,0) -- (101:1) arc (101:274:1) -- cycle;
  % Molt 274-360
  \fill[udnavy!60] (0,0) -- (274:1) arc (274:360:1) -- cycle;
  \draw[white, line width=0.8pt] (0,0) circle (1);
  \node[font=\scriptsize, white] at (14:0.7) {Gens};
  \node[font=\scriptsize, white] at (65:0.62) {Poc};
  \node[font=\scriptsize, white] at (187:0.6) {Bastant};
  \node[font=\scriptsize, white] at (317:0.64) {Molt};
\end{tikzpicture}

\figcap{«Com de segur et sents?»: Bastant $48\%$, Molt $24\%$, Poc $20\%$, Gens $8\%$.}
\end{center}

\noindent\textbf{Lectura:} sumant «Bastant» i «Molt», un $72\%$ se sent segur; però un
$28\%$ ($8\%+20\%$) se sent «Gens» o «Poc» segur: és la part que la intervenció hauria
d'atendre.
\end{exemple}

\subsection{Exercicis resolts}

\begin{exresolt}[completar el buidatge d'una pregunta]
A la pregunta «Quantes hores dorms les nits d'escola?» han respost $20$ alumnes:
$4$ diuen «Menys de $6$ h», $7$ «Entre $6$ i $7$», $6$ «Entre $7$ i $8$» i $3$ «Més
de $8$». Fes la taula amb percentatges i comprova que quadra.
\solucio
Calculem cada percentatge com $\dfrac{\text{recompte}}{20}\per 100$:

\begin{center}
\begin{tabular}{l c c}
\toprule
\textbf{Hores de son} & \textbf{Recompte} & \textbf{\%} \\
\midrule
Menys de $6$ & $4$ & $20\%$ \\
Entre $6$ i $7$ & $7$ & $35\%$ \\
Entre $7$ i $8$ & $6$ & $30\%$ \\
Més de $8$ & $3$ & $15\%$ \\
\midrule
\textbf{Total} & $20$ & $100\%$ \\
\bottomrule
\end{tabular}
\end{center}
\textbf{Comprovació:} $4+7+6+3=20$ i $20+35+30+15=100\%$.
\resposta{$20\%$, $35\%$, $30\%$, $15\%$; sumen $100\%$ i els recomptes, $20$.}
\end{exresolt}

\begin{exresolt}[interpretar el buidatge per a una memòria]
Amb les dades de l'exercici anterior, escriu dues o tres frases d'interpretació, com
si fossin per a la memòria del projecte.
\solucio
Una redacció possible: «De $20$ alumnes enquestats, un $20\%$ dorm menys de $6$ hores
les nits d'escola, una franja preocupant des del punt de vista del descans. La majoria
($65\%$) dorm entre $6$ i $8$ hores, i només un $15\%$ supera les $8$ hores. Es
recomana treballar la higiene del son, especialment amb el grup que dorm menys de $6$
hores.»
\resposta{Interpretació breu: destaca el $20\%$ que dorm poc, la majoria entre $6$ i
$8$ h, i proposa una actuació.}
\end{exresolt}

\subsection{Exercicis per fer}

\noindent\textit{Treballa a l'ordinador amb les respostes de la teva enquesta (o les
que et doni el professorat). Anota fórmules i resultats.}

\begin{exercicis}
\begin{exlist}
  \item Tens les respostes d'una pregunta a la columna A. Escriu la fórmula
        \texttt{=COMPTA.SI(...)} per comptar quantes vegades surt una opció concreta.
  \item Completa la taula de buidatge d'aquesta pregunta ($N=20$): «Sí» $8$, «No» $5$,
        «De vegades» $7$. Calcula el percentatge de cada opció.
  \item Comprova que els recomptes de l'exercici 2 sumen $N=20$ i que els percentatges
        sumen $100\%$.
  \item Calcula quants \textbf{graus} ocuparia cada opció de l'exercici 2 en un gràfic
        de sectors.
  \item \textbf{(Interpretació)} A partir del buidatge de l'exercici 2, escriu dues
        frases de conclusió per a una memòria, destacant què és el més rellevant.
  \item \textbf{(Raonament)} Per què és important \textbf{comprovar} que els
        percentatges sumen $100\%$ després de tabular? Quin tipus d'error detectaries
        si no sumessin exactament $100\%$?
\end{exlist}
\end{exercicis}

% ---------------------------------------------------------------------
\fullsolucions{Sessió 13}
\begin{solucions}
\begin{sollist}
  \item Per exemple, per comptar els «Sí» del rang \texttt{A2:A21}:
        \texttt{=COMPTA.SI(A2:A21;"Sí")} (en anglès \texttt{=COUNTIF(A2:A21,"Sí")}).
  \item «Sí»: $\dfrac{8}{20}=40\%$; «No»: $\dfrac{5}{20}=25\%$; «De vegades»:
        $\dfrac{7}{20}=35\%$.
  \item $8+5+7=20=N$: correcte. $40+25+35=100\%$: correcte.
  \item Graus ($\text{rel}\per 360^\circ$): «Sí» $0,40\per 360=144^\circ$; «No»
        $0,25\per 360=90^\circ$; «De vegades» $0,35\per 360=126^\circ$. (Comprovació:
        $144+90+126=360^\circ$.)
  \item Resposta oberta. Exemple: «Un $40\%$ de l'alumnat respon que sí i un $35\%$ que
        de vegades; només un $25\%$ respon que no. La majoria, doncs, es troba en la
        situació plantejada, almenys de manera ocasional.»
  \item Perquè si no sumen $100\%$ hi ha un \textbf{error}: algun recompte malament,
        una resposta no comptada, una resposta comptada dues vegades o categories que
        se solapen. La comprovació del $100\%$ és una manera ràpida de detectar errors
        de buidatge abans de presentar les dades.
\end{sollist}
\end{solucions}
